\clearpage
\section{Anexo 8.H: Reservas, perfiles y programación del esfuerzo}
\label{sec:cierre-reservas}
La exposición inicial de los 35 riesgos es de 1.183,11 HH, calculada con las probabilidades y los esfuerzos condicionales del Anexo 8.B. Los controles técnicos de diseño, innovación y calidad forman parte del trabajo ordinario; una diferencia de esfuerzo requiere estimación y aprobación específica.

La propuesta de contingencia específica es la suma p1 H0 de los 35 eventos: incluye atención contingente C y remanente L=H0-C. La prevención A es trabajo planificado adicional y se solicita como tal, separado de reserva. No se suman otra vez la exposición inicial, p1L y T=A+p1C a p1H0. La estimación corresponde al valor esperado; no determina percentiles ni suficiencia estadística de la reserva. Las hipótesis p1 dependen de evidencia; si no se verifica, se retiene la exposición inicial para replanificar y escalar, sin gastar una reserva proyectada como si ya estuviera aprobada.

La reserva protegida de S7, Tabla 7.4, pp. 45--46, es 3.072 HH: 1.152 en marcha blanca E1 (13--15) y 1.920 en producción/estabilización E1 (16--20). Ya integra 219.727,83 HH programadas sobre 216.655,83 HH base/cobertura. No se suma nuevamente a esos totales ni se presta a E2. El crédito inicial contra la contingencia específica es cero y sólo podrá reconocerse tras verificar la coincidencia de capacidades, perfiles, causas y períodos. JP podrá reconocer solapamiento sólo con la triple correspondencia riesgo--paquete--perfil/mes y acta de capacidad; al hacerlo, se resta exactamente una vez de la propuesta adicional, conservando la protección E1. Las fichas preventivas también se descuentan si ya están incluidas en paquetes.

La reserva de gestión propuesta es de 160 HH, sujeta a autorización del Comité Ejecutivo: cuatro ciclos extraordinarios por causas todavía no identificadas, uno por bloque 1--12, 13--20, 21--38 y 39--56. Cada ciclo reúne JP 8, ARQ 16, CAL 8 y SRE 8 HH: 4 por 40=160 HH. Son tolerancias de revisión, diagnóstico, coordinación y replanificación de un asunto emergente; no incluyen remediación de un riesgo conocido, controles técnicos ordinarios, cambios normativos ya registrados, correcciones ordinarias ni los 35 eventos. Si el asunto se identifica como un riesgo de T-16 se reclasifica, se anula su doble imputación y se usa su contingencia específica. Si excede el sobre de 40 HH, el Comité Ejecutivo decide capacidad/cambio antes de comprometer más trabajo. Sensibilidad de tolerancia: dos ciclos por bloque exigirían 320 HH, sin adoptar ese importe ni asegurar cobertura de eventos desconocidos.

La tabla mensual distribuye p1H0 uniformemente entre los meses de atención propuestos de cada ficha, 60\% al dueño y 40\% al apoyo (CAL, o ARQ si CAL es dueño). No predice fechas de ocurrencia ni cambia la ventana probabilística. Las cantidades de cálculo conservan precisión; el redondeo de impresión a 0,01 HH puede crear diferencias de centésimas al sumar celdas; los totales se calculan sobre los valores sin redondear. Prevención y gestión tienen sus propios calendarios y no están incluidas en esa tabla. La reserva distribuida es exposición esperada, no dotación comprometida: JP debe recalcular la ruta crítica y nivelar perfiles antes de solicitar su aprobación. No se convierten HH directamente a días.

\begingroup
\renewcommand{\thetable}{8.H.1}
\begin{landscape}
\setlength{\textwidth}{\linewidth}
\renewcommand{\arraystretch}{1.08}
\begin{tablalafrox}{Conciliación de esfuerzo y reservas, HH. Fuente: 8.B/8.F, S7 Tabla 7.4 y escenario explícito de gestión.}{tab:cierre-reserva}{F{0.28} G{0.14} Y}{\cab{Concepto} & \cab{HH} & \cab{Tratamiento y no duplicación}}
Exposición inicial vigente & 1.183,11 & Valoración previa al tratamiento; no sumar a la proyección posterior. \\
Trabajo preventivo adicional A & 180,00 & Trabajo planificado, 60/40 por ficha; descontar coincidencia con EDT antes de aprobar. \\
Contingencia específica propuesta p1H0 & 821,13 & Incluye C y L; condición: evidencia que sustenta p1. \\
Dentro de contingencia: p1C & 137,20 & Subconjunto; no reserva adicional. \\
Dentro de contingencia: p1L & 683,93 & Subconjunto de contingencia, no sumar otra vez. \\
Esfuerzo esperado posterior A+p1H0 & 1001,13 & Prevención más contingencia específica; sin gestión ni protección S7. \\
Gestión propuesta & 160,00 & Cuatro ciclos extraordinarios de 40 HH, con perfiles y bloques definidos. \\
Protección ya programada en S7 & 3.072,00 & Sólo E1; 1.152+1.920. Crédito demostrado inicialmente cero. \\
\end{tablalafrox}
\end{landscape}
\endgroup

\begingroup
\renewcommand{\thetable}{8.H.2}
\begin{landscape}
\setlength{\textwidth}{\linewidth}
\renewcommand{\arraystretch}{1.08}
\begin{tablalafrox}{Trabajo preventivo adicional propuesto por mes de primera atención y perfil (HH). Fuente: fichas 8.F.1; ejecutar antes del primer hito dependiente.}{tab:cierre-preventiva}{G{0.07} G{0.10} G{0.10} G{0.10} G{0.10} G{0.10} G{0.10} G{0.10}}{\cab{Mes} & \cab{JP} & \cab{ARQ} & \cab{DAT} & \cab{CAL} & \cab{SRE} & \cab{IMP} & \cab{SEG}}
1 & 8,40 & 18,00 & 22,80 & 48,40 & 2,40 & 15,60 & 2,40 \\
2 & 0,00 & 0,00 & 0,00 & 5,60 & 0,00 & 8,40 & 0,00 \\
13 & 0,00 & 8,40 & 0,00 & 13,60 & 0,00 & 12,00 & 0,00 \\
19 & 4,80 & 0,00 & 0,00 & 3,20 & 0,00 & 0,00 & 0,00 \\
20 & 3,60 & 0,00 & 0,00 & 2,40 & 0,00 & 0,00 & 0,00 \\
\end{tablalafrox}
\end{landscape}
\endgroup

\begingroup
\renewcommand{\thetable}{8.H.3}
\begin{landscape}
\setlength{\textwidth}{\linewidth}
\renewcommand{\arraystretch}{1.08}
\begin{tablalafrox}{Distribución mensual propuesta de contingencia específica (HH). Fuente: p1H0 y ventanas de 8.F.}{tab:cierre-mensual}{G{0.06} G{0.09} G{0.09} G{0.09} G{0.09} G{0.09} G{0.09} G{0.09} G{0.10}}{\cab{Mes} & \cab{JP} & \cab{ARQ} & \cab{DAT} & \cab{CAL} & \cab{SRE} & \cab{IMP} & \cab{SEG} & \cab{Total}}
1 & 1,01 & 6,86 & 7,32 & 12,49 & 0,13 & 2,98 & 0,05 & 30,83 \\
2 & 1,01 & 6,86 & 7,32 & 12,93 & 0,13 & 3,65 & 0,05 & 31,95 \\
3 & 1,01 & 2,70 & 7,32 & 10,16 & 0,13 & 3,65 & 0,05 & 25,02 \\
4 & 1,01 & 2,70 & 7,32 & 10,16 & 0,13 & 3,65 & 0,05 & 25,02 \\
5 & 1,01 & 2,70 & 7,32 & 10,16 & 0,13 & 3,65 & 0,05 & 25,02 \\
6 & 1,01 & 2,70 & 7,32 & 10,16 & 0,13 & 3,65 & 0,05 & 25,02 \\
7 & 1,01 & 2,70 & 7,32 & 10,16 & 0,13 & 3,65 & 0,05 & 25,02 \\
8 & 1,01 & 2,70 & 7,32 & 10,16 & 0,13 & 3,65 & 0,05 & 25,02 \\
9 & 1,01 & 2,70 & 7,32 & 10,16 & 0,13 & 3,65 & 0,05 & 25,02 \\
10 & 1,01 & 2,70 & 7,32 & 10,16 & 0,13 & 3,65 & 0,05 & 25,02 \\
11 & 1,01 & 2,70 & 7,32 & 10,16 & 0,13 & 3,65 & 0,05 & 25,02 \\
12 & 1,01 & 2,70 & 7,32 & 10,16 & 0,13 & 3,65 & 0,05 & 25,02 \\
13 & 1,01 & 5,02 & 7,32 & 12,76 & 0,19 & 5,17 & 0,05 & 31,52 \\
14 & 1,01 & 5,02 & 7,32 & 12,76 & 0,19 & 5,17 & 0,05 & 31,52 \\
15 & 1,01 & 5,02 & 7,32 & 12,76 & 0,19 & 5,17 & 0,05 & 31,52 \\
16 & 1,01 & 5,02 & 7,32 & 12,22 & 0,19 & 4,34 & 0,05 & 30,16 \\
17 & 1,01 & 5,02 & 7,32 & 12,22 & 0,19 & 4,34 & 0,05 & 30,16 \\
18 & 1,01 & 5,02 & 7,32 & 12,22 & 0,19 & 4,34 & 0,05 & 30,16 \\
19 & 8,71 & 3,39 & 5,36 & 14,96 & 0,19 & 4,34 & 0,05 & 37,00 \\
20 & 8,96 & 3,39 & 5,36 & 15,12 & 0,19 & 4,34 & 0,05 & 37,42 \\
21 & 0,75 & 1,13 & 0,00 & 2,76 & 0,19 & 2,03 & 0,05 & 6,91 \\
22 & 0,75 & 1,13 & 0,00 & 2,76 & 0,19 & 2,03 & 0,05 & 6,91 \\
23 & 0,75 & 1,13 & 0,00 & 2,76 & 0,19 & 2,03 & 0,05 & 6,91 \\
24 & 0,75 & 1,13 & 0,00 & 2,76 & 0,19 & 2,03 & 0,05 & 6,91 \\
25 & 0,75 & 1,13 & 0,00 & 2,76 & 0,19 & 2,03 & 0,05 & 6,91 \\
26 & 0,75 & 1,13 & 0,00 & 2,76 & 0,19 & 2,03 & 0,05 & 6,91 \\
27 & 0,75 & 1,13 & 0,00 & 2,76 & 0,19 & 2,03 & 0,05 & 6,91 \\
28 & 0,75 & 1,13 & 0,00 & 2,76 & 0,19 & 2,03 & 0,05 & 6,91 \\
29 & 0,75 & 1,13 & 0,00 & 2,76 & 0,19 & 2,03 & 0,05 & 6,91 \\
30 & 0,75 & 1,13 & 0,00 & 2,76 & 0,19 & 2,03 & 0,05 & 6,91 \\
31 & 0,75 & 1,13 & 0,00 & 2,76 & 0,19 & 2,03 & 0,05 & 6,91 \\
32 & 0,75 & 1,13 & 0,00 & 2,76 & 0,19 & 2,03 & 0,05 & 6,91 \\
33 & 0,75 & 1,13 & 0,00 & 2,76 & 0,19 & 2,03 & 0,05 & 6,91 \\
34 & 0,75 & 1,13 & 0,00 & 2,76 & 0,19 & 2,03 & 0,05 & 6,91 \\
35 & 0,75 & 1,13 & 0,00 & 2,76 & 0,19 & 2,03 & 0,05 & 6,91 \\
36 & 0,75 & 1,13 & 0,00 & 2,76 & 0,19 & 2,03 & 0,05 & 6,91 \\
37 & 0,75 & 1,13 & 0,00 & 2,76 & 0,19 & 2,03 & 0,05 & 6,91 \\
38 & 0,75 & 1,13 & 0,00 & 2,76 & 0,19 & 2,03 & 0,05 & 6,91 \\
39 & 0,75 & 1,13 & 0,00 & 2,76 & 0,19 & 2,03 & 0,05 & 6,91 \\
40 & 0,75 & 1,13 & 0,00 & 2,76 & 0,19 & 2,03 & 0,05 & 6,91 \\
41 & 0,75 & 1,13 & 0,00 & 2,76 & 0,19 & 2,03 & 0,05 & 6,91 \\
42 & 0,75 & 1,13 & 0,00 & 2,76 & 0,19 & 2,03 & 0,05 & 6,91 \\
43 & 0,75 & 1,13 & 0,00 & 2,76 & 0,19 & 2,03 & 0,05 & 6,91 \\
44 & 0,75 & 1,13 & 0,00 & 2,76 & 0,19 & 2,03 & 0,05 & 6,91 \\
45 & 0,75 & 1,13 & 0,00 & 2,76 & 0,19 & 2,03 & 0,05 & 6,91 \\
46 & 0,75 & 1,13 & 0,00 & 2,76 & 0,19 & 2,03 & 0,05 & 6,91 \\
47 & 0,75 & 1,13 & 0,00 & 2,76 & 0,19 & 2,03 & 0,05 & 6,91 \\
48 & 0,75 & 1,13 & 0,00 & 2,76 & 0,19 & 2,03 & 0,05 & 6,91 \\
49 & 0,75 & 1,13 & 0,00 & 2,76 & 0,19 & 2,03 & 0,05 & 6,91 \\
50 & 0,75 & 1,13 & 0,00 & 2,76 & 0,19 & 2,03 & 0,05 & 6,91 \\
51 & 0,75 & 1,13 & 0,00 & 2,76 & 0,19 & 2,03 & 0,05 & 6,91 \\
52 & 0,75 & 1,13 & 0,00 & 2,76 & 0,19 & 2,03 & 0,05 & 6,91 \\
53 & 0,75 & 1,13 & 0,00 & 2,76 & 0,19 & 2,03 & 0,05 & 6,91 \\
54 & 0,75 & 1,13 & 0,00 & 2,76 & 0,19 & 2,03 & 0,05 & 6,91 \\
55 & 0,75 & 1,13 & 0,00 & 2,76 & 0,19 & 2,03 & 0,05 & 6,91 \\
56 & 0,75 & 1,13 & 0,00 & 2,76 & 0,19 & 2,03 & 0,05 & 6,91 \\
\end{tablalafrox}
\end{landscape}
\endgroup
\paragraph{Programación y autorización.} JP asigna las fichas preventivas y la contingencia a los paquetes de migración/datos, implantación, integración, soporte, gobierno o salida que atienden cada evento. Antes de autorizar capacidad adicional debe verificar paquete, perfil, mes y capacidad existente, descontando toda coincidencia. La prevención se programa en el mes de primera atención y antes del hito dependiente; la contingencia se nivela con la distribución mensual propuesta. Las 3.072 HH protegidas mantienen su uso exclusivo en E1. Los ciclos de gestión se sitúan en los meses 1, 13, 21 y 39 como capacidad de decisión; sólo se consumen ante un asunto emergente autorizado. Las familias DAT/IMP conservan su clasificación de esfuerzo y SRE asume responsabilidad desde mes 22, sin duplicar cargas por relevo.

La programación está sujeta a nivelación y autorización de la línea base. JP verifica precedencias, holguras y ruta crítica, preservando congelamientos y las cuatro semanas finales de evidencia. Si la capacidad no permite cumplir, eleva refuerzo o cambio antes del hito, sin desplazarlo unilateralmente. La valoración económica utiliza la misma diferencia aprobada por mes y perfil, multiplicada por la tarifa neta validada. Gestión y protección no se registran como costo consumido por el solo hecho de reservarlas. JP presenta la conciliación al arranque, la actualiza en cada comité quincenal y antes de H5/H10; el Comité Ejecutivo autoriza los cambios y CAL verifica la ausencia de doble imputación. Una propuesta de reserva no acredita financiamiento aprobado.
